	\section{Summary: concentration normalisation of adsorption energy}

	\subsection{The problem: incomparable parameters across models}

	Different isotherm models parameterise the adsorbate--adsorbent
	interaction in different ways.  For example, the Langmuir constant
	$K_L$ has units $[C]^{-1}$, the Sips constant $K_S$ has units
	$[C]^{-n}$ (where $n$ depends on the heterogeneity), and the Temkin
	parameter $b_T$ has units of energy.
	Directly comparing $K_L = 0.05~\mathrm{L\,mg^{-1}}$ with
	$K_S = 0.12~\mathrm{(L\,mg^{-1})^{0.7}}$ is meaningless because
	the numbers refer to different physical quantities with different
	dimensions.  The same difficulty arises when comparing the ``best
	model'' parameters of two different adsorbents that happen to be
	described by different isotherm equations.

	A systematic solution is to extract, from every fitted model, a
	small set of \textbf{model-independent figures of merit} that
	(i)~have the same units regardless of the model chosen and
	(ii)~admit a clear thermodynamic or physical interpretation.

	\subsubsection{Three universal descriptors: $Q_{\max}$, $K_H$ and $K_{\mathrm{aff}}$}

	For any isotherm model with a well-defined saturation plateau and
	a finite initial slope, the following three quantities can always
	be extracted:

	\paragraph{Maximum adsorption capacity, $Q_{\max}$.}
	\begin{equation}
		Q_{\max} = \lim_{C_e \to \infty} q_e.
		\label{eq:Qmax_universal}
	\end{equation}
	Units: $[\text{mass}_{\text{adsorbate}} / \text{mass}_{\text{adsorbent}}]$
	(e.g.\ mg\,g$^{-1}$, mmol\,g$^{-1}$).
	This is the zeroth-order descriptor: \emph{how much} can the
	material adsorb.

	\paragraph{Henry constant, $K_H$.}
	\begin{equation}
		K_H = \left.\frac{\mathrm{d}q_e}{\mathrm{d}C_e}
		\right|_{C_e \to 0}.
		\label{eq:KH_universal}
	\end{equation}
	Units: $[q]/[C]$ (e.g.\ L\,g$^{-1}$).
	This is the first-order descriptor: \emph{how strongly} the
	material binds at infinite dilution.  $K_H$ controls the
	performance of single-use or trace-level applications.

	\paragraph{Affinity constant, $K_{\mathrm{aff}}$.}
	\begin{equation}
		\boxed{
			K_{\mathrm{aff}}
			= \frac{K_H}{Q_{\max}}
		}
		\label{eq:Kaff_universal}
	\end{equation}
	Units: $[C]^{-1}$ (e.g.\ L\,mg$^{-1}$, L\,mmol$^{-1}$).
	This is the capacity-normalised affinity: it measures the binding
	strength \emph{per unit of available capacity}.  Because it has
	units of inverse concentration, $K_{\mathrm{aff}}$ is the natural
	quantity to convert into a dimensionless equilibrium constant and
	thence into an adsorption free energy:
	\begin{equation}
		K_{\mathrm{aff}}^{\ast}
		= K_{\mathrm{aff}}\, C^\circ,
		\qquad
		\Delta G_{\mathrm{ads}}^{0}
		= -RT\ln K_{\mathrm{aff}}^{\ast}
		= -RT\ln(K_{\mathrm{aff}}\, C^\circ),
		\label{eq:DG_from_Kaff}
	\end{equation}
	where $C^\circ$ is a standard-state concentration.

	\subsubsection{Extraction of $Q_{\max}$, $K_H$ and $K_{\mathrm{aff}}$ from each model}

	Table~\ref{tab:universal_descriptors} shows how the three
	descriptors are obtained from the fitted parameters of each
	isotherm model.

	\begin{table}[htbp]
		\centering
		\small
		\caption{Universal descriptors extracted from each isotherm
			model.  $Q_{\max}$ is the saturation capacity,
			$K_H = \mathrm{d}q_e/\mathrm{d}C_e|_{C_e\to 0}$ is the
			Henry constant, and
			$K_{\mathrm{aff}} = K_H / Q_{\max}$ is the
			capacity-normalised affinity constant (units $[C]^{-1}$).
			Models marked with $\dagger$ have a divergent $Q_{\max}$
			or $K_H$, so only an effective (concentration-dependent)
			descriptor can be defined.}
		\label{tab:universal_descriptors}
		\renewcommand{\arraystretch}{1.4}
		\begin{tabular}{lccc}
			\hline
			\textbf{Model}
			& $\boldsymbol{Q_{\max}}$
			& $\boldsymbol{K_H}$
			& $\boldsymbol{K_{\mathrm{aff}} = K_H/Q_{\max}}$
			\\
			\hline
			Langmuir
			& $Q_{\max}$
			& $Q_{\max} K_L$
			& $K_L$
			\\
			Double Langmuir
			& $Q_1 + Q_2$
			& $Q_1 K_1 + Q_2 K_2$
			& $\dfrac{Q_1 K_1 + Q_2 K_2}{Q_1 + Q_2}$
			\\[4pt]
			Sips
			& $Q_{\max}$
			& $\infty$\;($n<1$)$^{\dagger}$
			& $K_S^{1/n}$\;(effective)
			\\
			T\'oth
			& $q_{\max}$
			& $q_{\max}\,b$
			& $b$
			\\
			Jovanovic
			& $q_{\max}$
			& $q_{\max}\,b$
			& $b$
			\\
			Temkin$^{\dagger}$
			& $\infty$
			& $\infty$
			& $a_T$\;(strongest site)
			\\
			UNILAN
			& $Q_{\max}$
			& $Q_{\max}\!\left\langle K \right\rangle_{\!f}$
			& $\sqrt{b_{\max}\,b_{\min}}$\;(geom.\ mean)
			\\
			Redlich--Peterson
			& $K_{RP}/a_{RP}$\;($\beta\!=\!1$)
			& $K_{RP}$
			& $a_{RP}$\;($\beta\!=\!1$)
			\\
			& $\infty$\;($\beta\!<\!1$)$^{\dagger}$
			&
			&
			\\
			Koble--Corrigan
			& $A_{KC}/B_{KC}$
			& $\infty$\;($1/n<1$)$^{\dagger}$
			& $B_{KC}^{\,n}$\;(effective)
			\\
			Freundlich$^{\dagger}$
			& $\infty$
			& $\infty$
			& $K_F^{\,n}\,(C^\circ)^{n-1}$\;(eff.)
			\\
			D-R / D-A
			& $Q_{\max}$
			& $\infty^{\dagger}$
			& ---\;(Polanyi framework)
			\\
			Hill
			& $Q_{\max}$
			& $\infty$\;($n_H\!<\!1$)$^{\dagger}$
			& $K_D^{-1/n_H}$\;(effective)
			\\
			\hline
		\end{tabular}
	\end{table}

	\paragraph{Notes on Table~\ref{tab:universal_descriptors}.}

	\begin{enumerate}
		\item For models with a finite Henry constant (Langmuir, T\'oth,
		Jovanovic, UNILAN, Redlich--Peterson with $\beta = 1$), the
		extraction is straightforward: $K_H$ is read from the low-$C_e$
		expansion and $K_{\mathrm{aff}} = K_H / Q_{\max}$.

		\item For models where the initial slope diverges (Sips with
		$n < 1$, Hill with $n_H < 1$, Freundlich, D-R/D-A), a true
		Henry constant does not exist.  In these cases an
		\emph{effective} affinity constant
		$K_{\mathrm{aff}}^{\mathrm{eff}}$ is defined by requiring
		that $K_{\mathrm{aff}}^{\mathrm{eff}}$ have units $[C]^{-1}$
		and reduce to the Langmuir result in the homogeneous limit.
		For the Sips model this gives
		$K_{\mathrm{aff}}^{\mathrm{eff}} = K_S^{1/n}$
		(Eq.~\eqref{eq:E_sips}).

		\item The Freundlich model is doubly degenerate: both
		$Q_{\max} = \infty$ and $K_H = \infty$.  An effective
		affinity can still be defined at a chosen reference
		concentration $C_{\mathrm{ref}}$ as
		$K_{\mathrm{aff}}^{\mathrm{eff}}(C_{\mathrm{ref}})
		= K_F\,(1/n)\,C_{\mathrm{ref}}^{1/n - 1}$,
		but the result is inherently concentration-dependent.
	\end{enumerate}

	\subsubsection{Concentration normalisation and the energy scale}

	Every isotherm model that admits a finite $K_{\mathrm{aff}}$ can be
	brought into a common energetic framework by the substitution
	\begin{equation}
		C_e = C^{\ast}\exp\!\left(-\frac{E}{RT}\right),
		\label{eq:conc_norm_general}
	\end{equation}
	where $C^{\ast}$ is a reference concentration and $E$ is the
	adsorption energy.  This \emph{concentration normalisation} maps
	each isotherm $q(C_e)$ into an energy-domain coverage $q(E)$,
	from which the site-energy distribution
	$f(E) = -\mathrm{d}q/\mathrm{d}E$ is obtained by
	differentiation.  The characteristic energy
	\begin{equation}
		E_c = RT\ln(K_{\mathrm{aff}}\,C^{\ast})
	\end{equation}
	locates the centre (or mode) of the distribution.

	Table~\ref{tab:energy_summary} collects the energy-distribution
	results for all models.

	\begin{table}[htbp]
		\centering
		\small
		\caption{Energy-distribution properties obtained via
			concentration normalisation
			$C_e = C^{\ast}\mathrm{e}^{-E/RT}$.
			$K_{\mathrm{aff}}$ is taken from
			Table~\ref{tab:universal_descriptors}.  The affinity
			energy is
			$E_0/RT = \ln(K_{\mathrm{aff}}\,C^\circ)$ in all cases.}
		\label{tab:energy_summary}
		\renewcommand{\arraystretch}{1.35}
		\begin{tabular}{lcccc}
			\hline
			\textbf{Model}
			& $\boldsymbol{K_{\mathrm{aff}}}$
			& \textbf{Distribution $f(E)$}
			& $\boldsymbol{\sigma_E}$
			& \textbf{Eq.}
			\\
			\hline
			Langmuir
			& $K_L$
			& Dirac $\delta(E - E_c)$
			& $0$
			& \eqref{eq:E_langmuir}
			\\
			Double Langmuir
			& $\dfrac{Q_1 K_1 + Q_2 K_2}{Q_1 + Q_2}$
			& Two Dirac deltas
			& ---
			& \eqref{eq:DeltaE12}
			\\
			Sips
			& $K_S^{1/n}$
			& Logistic
			& $\dfrac{\pi}{\sqrt{3}}\,nRT$
			& \eqref{eq:sigmaE_n_conc}
			\\
			T\'oth
			& $b$
			& Gen.\ logistic (type IV)
			& $> \dfrac{\pi}{\sqrt{3}}\,RT$\;($t<1$)
			& \eqref{eq:fE_toth}
			\\
			Jovanovic
			& $b$
			& Gumbel
			& $\dfrac{\pi}{\sqrt{6}}\,RT$
			& \eqref{eq:sigmaE_jov}
			\\
			Temkin
			& $a_T$
			& Uniform
			& $\dfrac{b_T Q_{\max}}{2\sqrt{3}}$
			& \eqref{eq:sigmaE_temkin}
			\\
			UNILAN
			& $\sqrt{b_{\max}b_{\min}}$
			& Uniform
			& $\dfrac{RT}{2\sqrt{3}}
			\ln\!\dfrac{b_{\max}}{b_{\min}}$
			& \eqref{eq:sigmaE_unilan}
			\\
			Redlich--Peterson
			& $K_{RP}$\,(Henry)
			& Non-standard
			& ---
			& \eqref{eq:E_RP_henry}
			\\
			& $a_{RP}^{1/\beta}$\,(high-$C$)
			&
			&
			& \eqref{eq:E_RP_high}
			\\
			Freundlich
			& $K_F^{\,n}(C^\circ)^{n-1}$
			& Power-law (unbounded)
			& $\infty$
			& ---
			\\
			D-R / D-A
			& ---
			& Gaussian-like
			& $E_a / \sqrt{2}$\,(D-R)
			& \eqref{eq:DA}
			\\
			Hill
			& $K_D^{-1/n_H}$
			& $\equiv$ Sips
			& $\dfrac{\pi}{\sqrt{3}}\,\dfrac{RT}{n_H}$
			& \eqref{eq:Hill}
			\\
			\hline
		\end{tabular}
	\end{table}

	\subsubsection{A universal figure of merit: the normalised affinity
		$K_{\mathrm{aff}}^{\ast}$}

	The central problem in comparing adsorbents is that the raw
	model parameters ($K_L$, $K_S$, $b$, $a_T$, \ldots) have
	different units, different physical meanings, and are tied to the
	particular isotherm model chosen.
	The quantity $K_{\mathrm{aff}} = K_H / Q_{\max}$ already solves the
	unit problem: it always has units $[C]^{-1}$.
	To obtain a \emph{dimensionless} figure of merit that can be
	compared across different studies, concentration units, and
	experimental conditions, we define the \textbf{normalised affinity}:
	\begin{equation}
		\boxed{
			K_{\mathrm{aff}}^{\ast}
			= K_{\mathrm{aff}} \times C^\circ
			= \frac{K_H}{Q_{\max}}\, C^\circ
		}
		\label{eq:Kaff_star}
	\end{equation}
	where $C^\circ$ is a \emph{declared} standard-state concentration.

	\paragraph{Choice of $C^\circ$.}
	The value of $C^\circ$ is conventional and must always be stated
	explicitly.  Two common choices are:
	\begin{itemize}
		\item $C^\circ = 1~\mathrm{mol\,L^{-1}}$
		(thermodynamic standard state for solutions).
		\item $C^\circ = 1~\mathrm{mg\,L^{-1}}$ (practical standard
		state for environmental and water-treatment studies, where
		concentrations are typically in ppm).
	\end{itemize}
	The choice affects the \emph{numerical value} of
	$K_{\mathrm{aff}}^{\ast}$ and the derived free energy, but
	\emph{not} the ranking of materials, provided the same $C^\circ$
	is used throughout a comparison.

	\paragraph{Physical interpretation.}
	$K_{\mathrm{aff}}^{\ast}$ is the equilibrium constant for the
	adsorption reaction at standard concentration, evaluated at the
	capacity-weighted mean site.
	\begin{itemize}
		\item $K_{\mathrm{aff}}^{\ast} \gg 1$: strong binding
		(essentially all solute captured at $C_e = C^\circ$).
		\item $K_{\mathrm{aff}}^{\ast} \approx 1$: moderate binding
		(half-saturation occurs near the standard concentration).
		\item $K_{\mathrm{aff}}^{\ast} \ll 1$: weak binding
		(negligible uptake at $C_e = C^\circ$).
	\end{itemize}

	\paragraph{Derived figures of merit.}
	From $K_{\mathrm{aff}}^{\ast}$, two additional descriptors follow
	immediately:
	\begin{align}
		\text{Affinity energy:}\qquad
		E_{\mathrm{aff}}
		&= RT\ln K_{\mathrm{aff}}^{\ast}
		= RT\ln(K_{\mathrm{aff}}\, C^\circ),
		\label{eq:Eaff_universal}
		\\[4pt]
		\text{Free energy of adsorption:}\qquad
		\Delta G_{\mathrm{ads}}^{0}
		&= -RT\ln K_{\mathrm{aff}}^{\ast}
		= -E_{\mathrm{aff}}.
		\label{eq:DG_universal}
	\end{align}
	These are the \emph{same} equations as the Langmuir free energy, but
	now applied to any model through the universal bridge
	$K_{\mathrm{aff}} = K_H / Q_{\max}$.

	\subsubsection{Practical protocol for cross-system comparison}

	The following step-by-step procedure makes reports from different
	laboratories, adsorbates, and models directly comparable.

	\begin{enumerate}
		\item \textbf{Declare the standard-state concentration
		$C^\circ$}~and the units of $C_e$ and $q_e$.
		Use $C^\circ = 1~\mathrm{mg\,L^{-1}}$ for environmental
		studies or $C^\circ = 1~\mathrm{mol\,L^{-1}}$ for
		thermodynamic consistency.

		\item \textbf{Fit the chosen isotherm model} by non-linear
		regression and obtain the native parameters
		($K_L$, $Q_{\max}$, $K_S$, $n$, $b$, $t$, etc.).

		\item \textbf{Extract the three universal descriptors}
		using Table~\ref{tab:universal_descriptors}:
		\begin{equation}
			Q_{\max},
			\qquad
			K_H,
			\qquad
			K_{\mathrm{aff}} = \frac{K_H}{Q_{\max}}.
			\label{eq:protocol_extract}
		\end{equation}
		If $K_H$ diverges (Sips, Freundlich, Hill with $n \neq 1$),
		use the effective $K_{\mathrm{aff}}$ from the table.

		\item \textbf{Compute the normalised affinity and energy:}
		\begin{equation}
			K_{\mathrm{aff}}^{\ast}
			= K_{\mathrm{aff}}\, C^\circ,
			\qquad
			E_{\mathrm{aff}} = RT\ln K_{\mathrm{aff}}^{\ast}.
		\end{equation}

		\item \textbf{Report the comparison set:}
		\begin{equation}
			\bigl\{
			Q_{\max},\;
			K_{\mathrm{aff}}^{\ast},\;
			E_{\mathrm{aff}},\;
			\sigma_E
			\bigr\}
			\label{eq:comparison_set}
		\end{equation}
		for every adsorbent--adsorbate pair.  These four numbers
		encode, respectively, capacity, affinity strength, affinity
		energy, and surface heterogeneity, and are comparable across
		models and laboratories.
	\end{enumerate}

	\paragraph{Example.}
	Consider an adsorbent fitted by a Sips model with
	$Q_{\max} = 120~\mathrm{mg\,g^{-1}}$,
	$K_S = 0.08~\mathrm{(L\,mg^{-1})^{0.7}}$, $n = 0.7$.
	A second adsorbent is described by a T\'oth model with
	$q_{\max} = 95~\mathrm{mg\,g^{-1}}$,
	$b = 0.15~\mathrm{L\,mg^{-1}}$, $t = 0.6$.
	Using $C^\circ = 1~\mathrm{mg\,L^{-1}}$:
	\begin{align}
		\text{Sips:}\quad
		K_{\mathrm{aff}}
		&= K_S^{1/n} = 0.08^{1/0.7}
		= 0.034~\mathrm{L\,mg^{-1}},
		\nonumber\\
		K_{\mathrm{aff}}^{\ast} &= 0.034,
		\qquad
		E_{\mathrm{aff}} / RT = \ln(0.034) = -3.38.
		\nonumber\\[4pt]
		\text{T\'oth:}\quad
		K_{\mathrm{aff}}
		&= b = 0.15~\mathrm{L\,mg^{-1}},
		\nonumber\\
		K_{\mathrm{aff}}^{\ast} &= 0.15,
		\qquad
		E_{\mathrm{aff}} / RT = \ln(0.15) = -1.90.
	\end{align}
	Despite using different models, the comparison is now
	straightforward: the T\'oth adsorbent has lower capacity
	($Q_{\max} = 95$ vs.\ 120~mg\,g$^{-1}$) but higher affinity
	($K_{\mathrm{aff}}^{\ast} = 0.15$ vs.\ 0.034), corresponding to
	a $1.48\,RT$ more favourable adsorption energy.

	\subsubsection{Why $K_{\mathrm{aff}}^{\ast}$ and not just
		$K_L$ or $K_S$?}

	It is common in the literature to report $\Delta G^0 = -RT\ln(K_L)$
	or $\Delta G^0 = -RT\ln(K_S)$ without specifying units or a
	standard-state concentration.  This practice leads to three
	systematic errors:

	\begin{enumerate}
		\item \textbf{Dimensional inconsistency.}
		If $K_L$ has units $\mathrm{L\,mg^{-1}}$ then
		$\ln(K_L)$ is undefined; the argument of a logarithm must be
		dimensionless.  Multiplying by $C^\circ$ cures this.

		\item \textbf{Unit dependence.}
		The same adsorbent gives different $\Delta G^0$ values
		depending on whether $C_e$ is expressed in
		$\mathrm{mg\,L^{-1}}$, $\mathrm{mmol\,L^{-1}}$, or
		$\mathrm{mol\,L^{-1}}$.
		The normalised affinity $K_{\mathrm{aff}}^{\ast}$, by
		construction, absorbs this ambiguity into the declared
		$C^\circ$.

		\item \textbf{Model dependence.}
		Comparing $K_L$ (Langmuir, units $[C]^{-1}$) with
		$K_S$ (Sips, units $[C]^{-n}$) is meaningless.
		The bridge $K_{\mathrm{aff}} = K_H / Q_{\max}$ reduces
		every model to a common $[C]^{-1}$ basis.
	\end{enumerate}

	\subsubsection{Summary of model-independent descriptors}

	Table~\ref{tab:comparison_protocol} summarises the recommended
	reporting protocol.

	\begin{table}[htbp]
		\centering
		\small
		\caption{Recommended model-independent descriptors for
			cross-system comparison of adsorption data.  All four
			quantities can be computed from any isotherm model via
			Table~\ref{tab:universal_descriptors}.}
		\label{tab:comparison_protocol}
		\renewcommand{\arraystretch}{1.4}
		\begin{tabular}{llll}
			\hline
			\textbf{Descriptor}
			& \textbf{Symbol}
			& \textbf{Units}
			& \textbf{Measures}
			\\
			\hline
			Saturation capacity
			& $Q_{\max}$
			& mg\,g$^{-1}$
			& How much
			\\
			Normalised affinity
			& $K_{\mathrm{aff}}^{\ast}
			= K_{\mathrm{aff}}\,C^\circ$
			& dimensionless
			& How strongly (per site)
			\\
			Affinity energy
			& $E_{\mathrm{aff}}
			= RT\ln K_{\mathrm{aff}}^{\ast}$
			& kJ\,mol$^{-1}$
			& Thermodynamic driving force
			\\
			Energy spread
			& $\sigma_E$
			& kJ\,mol$^{-1}$
			& Surface heterogeneity
			\\
			\hline
		\end{tabular}
	\end{table}

	\noindent\fbox{%
		\parbox{0.97\textwidth}{%
			\textbf{Key message.}
			Regardless of which isotherm model best fits the data, the
			four descriptors
			$\{Q_{\max},\, K_{\mathrm{aff}}^{\ast},\,
			E_{\mathrm{aff}},\, \sigma_E\}$
			provide a model-independent fingerprint of the
			adsorbent--adsorbate system.  Reporting these quantities
			--- together with the declared $C^\circ$ --- enables
			unambiguous cross-study comparison of adsorption
			performance.
		}%
	}

